\documentclass{article} % [paperfmt] 目标期刊要求 IEEEtran 或 acmart
\usepackage{graphicx}
\usepackage{amsmath, amssymb}
\usepackage{algorithm, algpseudocode}

\begin{document}

\title{A Lightweight Framework for Real-Time Object Detection in Edge Devices}
\author{Zhang Wei\thanks{Corresponding author. Email: wei.zhang@univ.edu} \and Li Ming \and Wang Fang} % [paperfmt] 双盲期未匿名
\maketitle

\begin{abstract}
Edge computing devices require efficient object detection models. We propose LightDet, a novel architecture that reduces computational cost by 40\% while maintaining accuracy. Experiments on COCO and custom datasets show state-of-the-art performance.
\end{abstract}

\section{Introduction}
Object detection has achieved remarkable success in recent years [1]. However, deploying these models on edge devices remains challenging due to limited computational resources. As shown in Fig. \ref{fig:arch}, our proposed pipeline addresses this gap.
In this paper, we make three main contributions: (1) a lightweight backbone, (2) an efficient neck module, and (3) a novel loss function. Eq. (1) demonstrates the core formulation.
\begin{equation}
\mathcal{L}_{total} = \lambda_1 \mathcal{L}_{cls} + \lambda_2 \mathcal{L}_{box} \label{eq:loss}
\end{equation}
We compare our method with recent baselines in Table 1. % [paperfmt] 硬编码 Table 1，未用 \ref

\begin{figure}[htbp]
\centering
\includegraphics[width=0.9\textwidth]{figs/pipeline.pdf}
\caption{The overall architecture of LightDet. It consists of three main stages.}
\label{fig:arch}
\end{figure}

\begin{table}[htbp]
\centering
\caption{Comparison with state-of-the-art methods on COCO val2017.} % [paperfmt] 若目标为IEEE，表注应在表上(正确)，但表格未用booktabs
\begin{tabular}{l c c c}
\hline
Method & mAP & Params (M) & FPS \\
\hline
YOLOv5s & 37.4 & 7.2 & 140 \\
YOLOv8n & 37.3 & 3.2 & 152 \\
Ours & \textbf{38.1} & 2.8 & \textbf{165} \\
\hline
\end{tabular}
\label{tab:results}
\end{table}

\section{Method}
Our approach builds upon the transformer-based paradigm. The attention mechanism is defined as:
\[
\text{Attention}(Q, K, V) = \text{softmax}\left(\frac{QK^T}{\sqrt{d_k}}\right)V
\]
To reduce redundancy, we introduce a spatial pruning module. As mentioned in Eq. (1), the loss function guides the optimization process. We implement the algorithm as follows:

\begin{algorithm}
\caption{LightDet Training Procedure}
\begin{algorithmic}[1]
\Require Dataset $\mathcal{D}$, epochs $E$, learning rate $\eta$
\Ensure Trained model weights $\theta$
\State Initialize $\theta$ with pretrained backbone
\For{epoch $e = 1$ to $E$}
    \State Sample batch $\mathcal{B} \sim \mathcal{D}$
    \State Compute $\mathcal{L}_{total}$ via Eq. (1)
    \State Update $\theta \leftarrow \theta - \eta \nabla_\theta \mathcal{L}$
\EndFor
\end{algorithmic}
\end{algorithm}

\section{Experiments}
We evaluate on the MS COCO dataset. All experiments use PyTorch 2.0 on a single RTX 4090. Table 1 summarizes the main results. Fig. \ref{fig:ablation} shows the ablation study.
\begin{figure}[htbp]
\centering
\includegraphics[width=0.85\textwidth]{figs/ablation.png} % [paperfmt] .png 可能 DPI 不足
\caption{Ablation study on different components.}
\label{fig:ablation}
\end{figure}

\section{Conclusion}
We presented LightDet, a fast and accurate detector for edge devices. Future work will focus on quantization-aware training.

\bibliographystyle{plain}
\bibliography{references}
\end{document}